\documentclass{article}
\usepackage{amsmath}
\usepackage{amssymb}
\usepackage{booktabs}
\usepackage{array}

\begin{document}

\title{�� Lecture Scribe \\ Marginal PMF, Correlation, Covariance, Joint Gaussian Random Variables \\ (Based strictly on Lecture-18 slides)}
\maketitle

\section{�� 1. Marginal PMF}

\subsection{✅ Definition}
For discrete random variables $X$ and $Y$, the joint PMF is \\
$p_{XY}(x,y)=P(X=x,Y=y)$ \\
The marginal PMF is obtained by summing the joint PMF over all possible values of the other variable. \\
Thus, \\
$p_X(x)=\sum_{y} p_{XY}(x,y)$ \quad $p_Y(y)=\sum_{x} p_{XY}(x,y)$ \\
\textbf{�� Key Idea:} \\
Marginal PMF of one variable = sum of joint PMF across all values of the other variable. \\
(Explained in slide definition)

\subsection{✅ Marginal PMF from Joint PMF Table}
From the joint table: \\
Row sums $\rightarrow$ give $p_X(x)$ \\
Column sums $\rightarrow$ give $p_Y(y)$ \\
Mathematically, \\
$p_X(x_i)=\sum_j p_{ij}$ \quad $p_Y(y_j)=\sum_i p_{ij}$ \\
Also, \\
$\sum_i p_X(x_i)=1,\quad \sum_j p_Y(y_j)=1$ \\
(as marginal PMFs must sum to 1).

\subsection{✅ Solved Example}
Joint PMF: \\
\begin{tabular}{c|cc}
$p_{XY}(x,y)$ & $y=0$ & $y=1$ \\
\hline
$x=0$ & $1/8$ & $1/4$ \\
$x=1$ & $3/8$ & $1/4$ \\
\end{tabular}

\textbf{Step-1: Find $p_X(x)$} \\
$p_X(0)=\frac{1}{8}+\frac{1}{4} =\frac{1}{8}+\frac{2}{8} =\frac{3}{8}$ \quad $p_X(1)=\frac{3}{8}+\frac{1}{4} =\frac{3}{8}+\frac{2}{8} =\frac{5}{8}$

\textbf{Step-2: Find $p_Y(y)$} \\
$p_Y(0)=\frac{1}{8}+\frac{3}{8} =\frac{4}{8} =\frac{1}{2}$ \quad $p_Y(1)=\frac{1}{4}+\frac{1}{4} =\frac{1}{2}$ \\
Thus, \\
$p_X(x)= \begin{cases} 3/8,&x=0\\ 5/8,&x=1 \end{cases}$ \quad $p_Y(y)= \begin{cases} 1/2,&y=0\\ 1/2,&y=1 \end{cases}$ \\
(Example derivation shown in slides)

\section{�� 2. Correlation}

\subsection{✅ Product Moment}
A basic measure of joint behaviour is \\
$R_{XY}=E[XY]$ \\
It represents the average value of the product $XY$. \\
\textbf{Interpretation} \\
Depends on scale of $X$ and $Y$ \\
Computed about origin (not about mean) \\
Mixes: \\
mean \\
spread \\
dependence \\
If \\
$R_{XY}=0$ \\
$\rightarrow$ variables are orthogonal about origin \\
To remove mean effect $\rightarrow$ center variables $\rightarrow$ covariance.

\section{�� 3. Covariance}

\subsection{✅ Definition}
$\text{Cov}(X,Y)=E[(X-\mu_X)(Y-\mu_Y)]$ \\
Where \\
$\mu_X=E[X],\quad \mu_Y=E[Y]$ \\
\textbf{Interpretation} \\
Covariance measures whether deviations from mean occur: \\
Same direction $\rightarrow$ Positive covariance \\
Opposite direction $\rightarrow$ Negative covariance \\
No linear relation $\rightarrow$ Zero covariance \\
Magnitude depends on units $\rightarrow$ cannot measure strength. \\
(Interpretation slide)

\subsection{✅ Algebraic Form (Step-by-Step Derivation)}
Start with: \\
$\text{Cov}(X,Y)=E[(X-\mu_X)(Y-\mu_Y)]$ \\
Expand: \\
$=E[XY-X\mu_Y-\mu_XY+\mu_X\mu_Y]$ \\
Apply expectation: \\
$=E[XY]-\mu_YE[X]-\mu_XE[Y]+\mu_X\mu_Y$ \\
Since: \\
$E[X]=\mu_X,\quad E[Y]=\mu_Y$ \\
Therefore: \\
$\text{Cov}(X,Y)=E[XY]-\mu_X\mu_Y$ \\
(Full derivation from slide)

\subsection{✅ Properties}
\textbf{Symmetry} \\
$\text{Cov}(X,Y)=\text{Cov}(Y,X)$ \\
\textbf{Variance special case} \\
$\text{Cov}(X,X)=Var(X)$ \\
\textbf{Linearity} \\
$\text{Cov}(aX+b,cY+d)=ac\,\text{Cov}(X,Y)$ \\
\textbf{Independence} \\
$X\perp Y \Rightarrow \text{Cov}(X,Y)=0$ \\
But \\
\textbf{�� Zero covariance $\neq$ independence (general case).} \\
(Properties slide)

\subsection{✅ Solved Example-1}
Given: \\
$Y=-2X+3$ \\
Find $\text{Cov}(X,Y)$ \\
Using: \\
$\text{Cov}(X,Y)=E[XY]-E[X]E[Y]$

\textbf{Step-1: Compute $XY$} \\
$XY=X(-2X+3)=-2X^2+3X$

\textbf{Step-2: Compute $E[Y]$} \\
$E[Y]=E[-2X+3]=-2E[X]+3$

\textbf{Step-3: Substitute} \\
$\text{Cov}(X,Y)=E[-2X^2+3X]-E[X](-2E[X]+3)$ \\
$=-2E[X^2]+3E[X]+2(E[X])^2-3E[X]$ \\
$=-2E[X^2]+2(E[X])^2$ \\
Since \\
$Var(X)=E[X^2]-(E[X])^2$ \\
Thus, \\
$\text{Cov}(X,Y)=-2Var(X)$ \\
(Example slide derivation)

\subsection{✅ Pearson Correlation Coefficient}
$\rho_{XY}=\frac{Cov(X,Y)}{\sqrt{Var(X)Var(Y)}} =\frac{Cov(X,Y)}{\sigma_X\sigma_Y}$ \\
\textbf{Interpretation} \\
Measures strength + direction of linear relationship \\
Normalized covariance \\
Unit-free \\
\textbf{Range:} \\
$-1\le \rho_{XY}\le 1$ \\
$>0 \rightarrow$ positive trend \\
$<0 \rightarrow$ negative trend \\
$=0 \rightarrow$ no linear correlation \\
(Slides explanation)

\subsection{✅ Solved Example-2 (Continuous Case)}
Joint PDF: \\
$f_{XY}(x,y)=\frac{xy}{96},\quad 0<x<4,\ 1<y<5$

\textbf{Step-1: Compute $E[X]$} \\
$E[X]=\int_0^4 x f_X(x)dx =\frac{1}{8}\int_0^4 x^2dx$ \\
$=\frac{1}{8}\left[\frac{x^3}{3}\right]_0^4 =\frac{1}{8}\cdot\frac{64}{3} =\frac{8}{3}$

\textbf{Step-2: Compute $E[Y]$} \\
$E[Y]=\int_1^5 y f_Y(y)dy =\frac{1}{12}\int_1^5 y^2dy$ \\
$=\frac{1}{12}\left[\frac{y^3}{3}\right]_1^5 =\frac{31}{9}$

\textbf{Step-3: Compute $E[XY]$} \\
Since $X,Y$ independent: \\
$E[XY]=E[X]E[Y] =\frac{8}{3}\cdot\frac{31}{9} =\frac{248}{27}$

\textbf{Step-4: Linear Expectation} \\
$E[2X+3Y]=2E[X]+3E[Y] =\frac{16}{3}+\frac{31}{3} =\frac{47}{3}$

\textbf{Step-5: Variance of $X$} \\
$E[X^2]=\frac{1}{8}\int_0^4 x^3dx=8$ \\
$Var(X)=8-\left(\frac{8}{3}\right)^2 =\frac{8}{9}$

\textbf{Step-6: Variance of $Y$} \\
$E[Y^2]=13$ \\
$Var(Y)=13-\left(\frac{31}{9}\right)^2 =\frac{92}{81}$

\textbf{Step-7: Covariance} \\
$\text{Cov}(X,Y)=E[XY]-E[X]E[Y] =\frac{248}{27}-\frac{248}{27}=0$

\textbf{Step-8: Correlation} \\
$\rho_{XY}=\frac{0}{\sqrt{Var(X)Var(Y)}}=0$ \\
(All steps from solution slides)

\section{�� 4. Joint Gaussian Random Variables}

\subsection{✅ Definition}
A pair $(X,Y)$ is jointly Gaussian if: \\
Joint density has bivariate Gaussian form \\
Marginals are Gaussian \\
\textbf{Applications:} \\
Signal processing \\
Communication systems \\
Noisy measurements modelling \\
(Definition slide)

\subsection{✅ Algebraic Form}
$f_{XY}(x,y)= \frac{1}{2\pi\sigma_X\sigma_Y\sqrt{1-\rho^2}} \exp\left( -\frac{1}{2(1-\rho^2)} \left[ \left(\frac{x-\mu_X}{\sigma_X}\right)^2 +\left(\frac{y-\mu_Y}{\sigma_Y}\right)^2 -2\rho\left(\frac{x-\mu_X}{\sigma_X}\right) \left(\frac{y-\mu_Y}{\sigma_Y}\right) \right] \right)$ \\
Where: \\
$\mu_X,\mu_Y \rightarrow$ means \\
$\sigma_X,\sigma_Y \rightarrow$ standard deviations \\
$\rho \rightarrow$ correlation parameter

\subsection{✅ Interpretation Example}
Example given: \\
$X$: CO₂ concentration \\
$Y$: Global temperature anomaly \\
If: \\
$\rho>0$ \\
$\rightarrow$ High CO₂ tends to occur with higher temperature. \\
(Interpretation slide)

\end{document}
